\documentclass[12pt]{article}
\usepackage[margin=0.75in]{geometry}

\newcommand{\duedate}{04/01/2026}
\newcommand{\assignment}{Online Meta-Learning: Summary}

\newcommand{\name}{Aksel Kretsinger-Walters, adk2164}
\newcommand{\email}{adk2164@columbia.edu}

\makeatletter
\def\input@path{{../}{../../}{../../../}}
\makeatother
\input{pset_template_CLMM.tex} %% DO NOT CHANGE THIS LINE

\begin{document}
\psetheader %% DO NOT CHANGE THIS LINE

\section*{Online Meta-Learning}

\paragraph{Motivation.}
Two distinct research paradigms address how agents can leverage past experience
to improve future learning. Meta-learning studies how to learn a prior over model
parameters for fast adaptation, but assumes tasks are available as a batch.
Online (regret-based) learning handles sequential, non-stationary task streams,
but trains a single model without task-specific adaptation.

Neither setting alone captures continual lifelong learning: meta-learning neglects
sequentiality, and online learning does not consider how past experience can
accelerate adaptation to new tasks.

\paragraph{Core Idea.}
The paper introduces the \textbf{online meta-learning} problem setting, which
merges ideas from both paradigms. An agent faces tasks one after another in a
sequential fashion and must simultaneously:

\begin{itemize}
		\item Use past experience to learn good priors (meta-learning).
		\item Adapt quickly to the current task at hand (online learning).
\end{itemize}

At each round $t$, the agent maintains parameters $\mathbf{w}_t$, the world
reveals a task defined by loss $f_t$, and the agent applies a local update
procedure $U_t$ (e.g., a gradient descent step) before being evaluated.

\paragraph{Follow the Meta Leader (FTML).}
The proposed algorithm extends MAML to the online setting by applying the
Follow the Leader (FTL) principle at the meta-level. FTML selects parameters
at each round by solving:
\[
		\mathbf{w}_{t+1} = \arg\min_{\mathbf{w}} \sum_{k=1}^{t} f_k\bigl(U_k(\mathbf{w})\bigr).
\]

This can be interpreted as playing the best meta-learner in hindsight up to
round $t$. The comparator class is powerful: it allows task-specific adaptation
via $U_t$, so achieving sublinear regret means the algorithm is competitive with
the best MAML-style initialization in hindsight.

\paragraph{Theoretical Guarantees.}
Under standard smoothness assumptions plus one additional condition on Lipschitz
Hessians ($C^2$-smoothness), the authors prove:

\begin{enumerate}
		\item The MAML-like objective inherits convexity from the original loss
				when the step size $\alpha$ is sufficiently small (Theorem 1).
		\item FTML achieves $O\!\left(\frac{G^2}{\mu}\log T\right)$ regret against
				the best meta-learner in hindsight (Corollary 2), providing the first
				provable and efficient guarantees for MAML-like objectives.
\end{enumerate}

\paragraph{Practical Algorithm.}
For deep networks with non-convex losses, the paper adapts FTML using stochastic
optimization. At each round, the algorithm maintains a task buffer $B$ and
computes meta-gradients by sampling past tasks and using independently drawn
minibatches for the inner update ($\mathcal{D}^{\text{tr}}_k$) and outer
optimization ($\mathcal{D}^{\text{val}}_k$) to reduce bias.

The gradient computation follows:
\[
		\mathbf{g}_t = \nabla_{\mathbf{w}} \, \mathbb{E}_{k \sim \nu^t}
		\mathcal{L}\bigl(\mathcal{D}_k^{\text{val}}, U_k(\mathbf{w})\bigr),
		\quad \text{where} \quad
		U_k(\mathbf{w}) \equiv \mathbf{w} - \alpha \, \nabla_{\mathbf{w}}
		\mathcal{L}(\mathcal{D}_k^{\text{tr}}, \mathbf{w}).
\]

\paragraph{Experimental Results.}
FTML is evaluated on three vision-based sequential learning problems:

\begin{itemize}
		\item \textbf{Rainbow MNIST:} 56 tasks with varying backgrounds, scales,
				and rotations. FTML learns new tasks faster over time and outperforms
				all baselines in both efficiency and final performance.
		\item \textbf{Five-Way CIFAR-100:} A sequence of 5-way classification tasks.
				FTML learns each task more efficiently than training from scratch and
				benefits from adapting all layers rather than only the last layer.
		\item \textbf{PASCAL 3D+ Pose Prediction:} 90 pose regression tasks. FTML
				solves many tasks with only 10 datapoints, demonstrating strong
				forward transfer.
\end{itemize}

Across all experiments, FTML substantially outperforms baselines including
Train on Everything (TOE), training from scratch, and FTL with fine-tuning.

\paragraph{Limitations.}
FTML inherits the computational cost of FTL, which grows as loss functions
accumulate. Storing all datapoints from prior tasks may also be impractical
at scale. The theoretical analysis is restricted to convex losses with
$C^2$-smoothness, leaving non-convex guarantees as future work.

\paragraph{Takeaway.}
Online Meta-Learning bridges meta-learning and online learning into a unified
framework for continual lifelong learning. The FTML algorithm demonstrates that
meta-learning sequentially---becoming more proficient at learning new tasks over
time---is both theoretically grounded and empirically effective.

\end{document}
